% ============================================================
%  ESTUDIOS CONTRASTADOS DIGESTIVOS
%  Licenciatura en Imagenología --- UdelaR
%  Dr. Nicolás Asteggiante -- Asis. Lic. Gonzalo Ferreira
%  Unidad Académica de Imagenología, 2026
% ============================================================

\chapter{Estudios contrastados digestivos}



% ============================================================
%  PUNTOS CLAVE PARA EL EXAMEN ORAL
% ============================================================
\begin{cajakeypoints}
	\textbf{\color{azulprincipal} Puntos clave para el examen oral}
	\begin{itemize}[noitemsep, topsep=2pt]
		\item Un estudio contrastado modifica la \textbf{densidad} del órgano mediante un fármaco (medio de contraste) y radiaciones ionizantes; tiene utilidad \textbf{anatómica y funcional}.
		\item Contrastes \textbf{positivos}: sulfato de bario (Z=56) y yodado hidrosoluble (Z=53). Contrastes \textbf{negativos}: aire, sales efervescentes. El \textbf{doble contraste} combina ambos.
		\item El bario se usa en estudios de rutina; el \textbf{yodado hidrosoluble} es obligatorio cuando se sospecha \textbf{perforación, fuga o fístula} (el bario en mediastino/peritoneo causa reacción inflamatoria grave).
		\item El \textbf{estudio de la deglución} evalúa disfagia orgánica o funcional; se inicia siempre con consistencia \textbf{semisólida} para evitar aspiración.
		\item En el \textbf{EGD} la posición de elección es la \textbf{OAD} (oblicua anterior derecha) para despejar el esófago de la columna. El paciente debe estar en \textbf{ayunas 6 horas} y no fumar/mascar chicle.
		\item El \textbf{colon por enema} se realiza con bario + aire (doble contraste); requiere preparación con laxante el día previo y cánula por vía anal en decúbito prono.
		\item La \textbf{colangiografía postoperatoria} usa contraste yodado diluido al 50\% inyectado por el tubo de Kehr; evalúa litiasis residual, estenosis y permeabilidad al duodeno.
		\item Signo patognomónico de \textbf{acalasia}: ``punta de lápiz'' o ``cola de ratón''. Signo de estenosis neoplásica: bordes ``en corazón de manzana'', rigidez segmentaria, defecto de llenado.
	\end{itemize}
\end{cajakeypoints}

\vspace{1em}

% ============================================================
%  SECCIÓN 1 --- DEFINICIÓN
% ============================================================
\section*{Definición}

\begin{cajadefinicion}
	Un \textbf{estudio contrastado} es un procedimiento médico que utiliza \textbf{radiaciones ionizantes} y la administración de un \textbf{medio de contraste} (fármaco de alto número atómico) para visualizar estructuras que, por sí solas, son indistinguibles de su entorno por tener densidad similar a los tejidos adyacentes.
\end{cajadefinicion}

Su objetivo es cambiar la densidad de órganos huecos (como esófago, estómago, colon o vía biliar) permitiendo su diferenciación en radiología convencional. Son estudios complementarios, solicitados habitualmente por el cirujano en contexto postoperatorio o por el gastroenterólogo. Su utilidad es \textbf{doble}: anatómica (forma, estructura, calibre) y funcional (motilidad, tránsito, competencia de esfínteres).

% ============================================================
%  SECCIÓN 2 --- CLASIFICACIÓN DE MEDIOS DE CONTRASTE
% ============================================================
\section*{Clasificación de los Medios de Contraste}

\subsection*{Contrastes positivos (aumentan densidad)}

{\renewcommand{\arraystretch}{1.35}
	\begin{center}
		\begin{tabular}{>{\bfseries}p{3.5cm} p{5.5cm} p{4.5cm}}
			\hline
			Contraste & Características & Uso principal \\
			\hline
			Sulfato de bario & Sal mineral inerte, viscosa; no se absorbe; se elimina por vía rectal. Se adhiere a la mucosa (\textit{mucosografía}). & Estudios de rutina del tubo digestivo: EGD, colon por enema, tránsito del delgado. \\[0.5em]
			\hline
			Yodado hidrosoluble & No iónico, baja osmolaridad (ej. Omnipaque, Z=53). No se adhiere a la mucosa. & Postoperatorios, sospecha de perforación/fuga/fístula, colangiografía, prequirúrgico de reconstrucción. \\
			\hline
		\end{tabular}
	\end{center}
}
.\\
\begin{cajakeypoints}
	\textbf{Regla de oro:} Si hay sospecha de que el contraste puede salir del tubo digestivo (perforación, anastomosis reciente, fístula), usar \textbf{siempre yodado hidrosoluble}. El bario libre en mediastino o peritoneo produce \textbf{mediastinitis o peritonitis química} de manejo muy difícil.
\end{cajakeypoints}

\subsection*{Contrastes negativos (disminuyen densidad)}

Se utilizan para distender el órgano y generar contraste mediante gas. En el \textbf{colon por enema} se usa \textbf{aire}; en el \textbf{EGD} se utilizan \textbf{sales efervescentes} (bicarbonato de sodio, carbonato de sodio, ácido cítrico; marca comercial tipo Uvasal). La combinación de un contraste positivo con uno negativo produce el \textbf{doble contraste}: el gas distiende la luz y el bario pinta la mucosa, permitiendo ver el detalle fino de los pliegues.

% ============================================================
%  SECCIÓN 3 --- ESTUDIOS DEL TUBO DIGESTIVO
% ============================================================
\section*{Estudios del Tubo Digestivo}

\subsection*{Estudio de la Deglución}

Evalúa la \textbf{disfagia} (dificultad para tragar), que puede ser \textbf{orgánica} (causa estructural: tumor, divertículo de Zenker, lesión cáustica, osteofitos) o \textbf{funcional} (causa neuromuscular: ACV, ELA, Alzheimer, miastenia gravis). Es un estudio \textbf{dinámico} realizado bajo \textbf{radioscopia en tiempo real}.

\textbf{Preparación:} Paciente en ayunas. Si tiene SNG, retirarla 12 horas antes.

\textbf{Materiales -- protocolo artesanal} (sulfato de bario en tres consistencias):
\begin{itemize}[noitemsep]
	\item \textbf{Líquida:} SB 170 mg en 100 cc de agua.
	\item \textbf{Semisólida:} SB 170 mg en 50 cc de agua + 5 mg de espesante.
	\item \textbf{Sólida:} Trozo de pan de 5 g sumergido en SB.
\end{itemize}

\textbf{Orden de administración (crítico para seguridad):} Siempre comenzar por la \textbf{consistencia semisólida}. Iniciar con líquido podría causar una aspiración masiva que impediría completar el estudio.

\textbf{Técnica:} El paciente realiza un buche sin tragar; bajo radioscopia el médico da la orden de deglutir y se captura una secuencia de disparos seriada (video) en \textbf{perfil} (tiempos de deglución) y \textbf{frente} (residuos en valéculas y senos piriformes). En un estudio normal no debe quedar residuo en estos sectores.

\textbf{Hallazgos críticos:}
\begin{itemize}[noitemsep]
	\item \textbf{Penetración:} el contraste entra al vestíbulo laríngeo pero se detiene en las cuerdas vocales.
	\item \textbf{Aspiración:} el contraste sobrepasa la glotis; se visualizan tráquea y bronquios ``pintados'' de bario.
\end{itemize}

\textbf{Rol del Licenciado:} Adquirir la secuencia de video completa. El médico asiste al paciente (frecuentemente en silla de ruedas o con movilidad reducida) mientras el licenciado opera el equipo.

% ---
\subsection*{Esofagogastroduodeno (EGD / SEGD)}

Evaluación anatómica y funcional del esófago, estómago y duodeno mediante doble contraste. Indicaciones principales: disfagia funcional u orgánica, ERGE, hernia hiatal, estenosis, postoperatorio de cirugías esofágicas o gástricas (cirugías de Nissen, bariátricas).

\textbf{Preparación:} Ayuno mínimo de 6 horas. No fumar, mascar chicle, ni tomar mate (producen saliva y gas que arruinan la adherencia del bario). No se usa Buscapina (para conservar la evaluación de la motilidad).

\textbf{Materiales:} Sales efervescentes disueltas en poca agua (el paciente las bebe rápido sin eructar) + bario diluido. Aproximadamente \textbf{5 chasis 35$\times$43}, apaisados y verticales según la fase.

\textbf{Parámetros técnicos:} Alto kVp (escala larga de grises para compensar el contraste del bario) y tiempos de exposición cortos (evitar borrosidad por peristaltismo). Referencia: 70 kV / 10--20 mAs según el sector.

\clearpage

\textbf{Secuencia de posiciones y fases:}

{\renewcommand{\arraystretch}{1.35}
	\begin{center}
		\begin{tabular}{>{\bfseries}p{3.5cm} p{5cm} p{5cm}}
			\hline
			\textbf{Posición} & \textbf{Dinámica del contraste} & \textbf{Objetivo} \\
			\hline
			Bipedestación perfil & Buche de bario diluido; evaluación dinámica de la faringe. & Descartar aspiración antes de proceder. \\[0.5em]
			\hline
			Bipedestación OAD & Paciente apoya hombro izquierdo, eleva lado derecho. & Despejar esófago de columna y corazón; esófago en repleción y doble contraste. \\[0.5em]
			\hline
			Bipedestación frente & Rayo horizontal. Nivel hidroaéreo visible. & Fundus en doble contraste (aire); cuerpo y antropíloro en repleción por gravedad. \\[0.5em]
			\hline
			Decúbito supino frente & Bario se desplaza a zonas posteriores. & Fundus y antro distal en repleción; cuerpo gástrico en doble contraste (anterior). \\[0.5em]
			\hline
			Decúbito OAD & Paciente gira hacia su izquierda. & Fundus en repleción; cuerpo y antropíloro en doble contraste. \\[0.5em]
			\hline
			Decúbito OAI + maniobras reflujo & Paciente gira hacia su derecha; Trendelenburg; deglución de saliva. & Fundus, antropíloro y bulbo en repleción; evaluación del esfínter esofágico inferior. \\[0.5em]
			\hline
			Decúbito prono / OPI & Casi boca abajo; bario pasa al duodeno por gravedad. & Documentar las 4 porciones del duodeno (bulbo, descendente, horizontal, ascendente). \\[0.5em]
			\hline
			Decúbito OAD final & Se vuelve a OAD tras movilizar al paciente. & Duodeno en doble contraste (la posición más difícil de lograr). \\
			\hline
		\end{tabular}
	\end{center}
}

\begin{cajakeypoints}
	\textbf{Anatomía del duodeno:} Bulbo (1\textordfeminine\ porción): liso, sin estriaciones. Porciones 2, 3 y 4: pliegues circulares. La 2\textordfeminine\ porción recibe los conductos biliar y pancreático. La 4\textordfeminine\ finaliza en el \textbf{ángulo de Treitz}.
\end{cajakeypoints}

\textbf{Maniobras de mesa:}
\begin{itemize}[noitemsep]
	\item \textbf{Trendelenburg} (cabeza abajo): rellena sectores superiores con bario; fuerza el reflujo gastroesofágico.
	\item \textbf{Anti-Trendelenburg} (cabeza arriba): desplaza bario hacia sectores inferiores; útil para doble contraste superior.
\end{itemize}


\textbf{Postoperatorio urgente:} Si hay sospecha de fuga antes del alta, usar \textbf{yodado hidrosoluble diluido} en lugar de bario. Control estándar postquirúrgico: a los 3 meses.

\clearpage
% ---
\subsection*{Tránsito Esofágico (postoperatorio de esofagectomía)}

Se indica específicamente para control tras \textbf{esofagectomía con ascenso gástrico} (resección de esófago por cáncer y reconstrucción con estómago). El punto crítico es la \textbf{anastomosis}. Se utiliza \textbf{yodado hidrosoluble} (riesgo de fuga al mediastino).

\textbf{Proyecciones obligatorias:} frente estricto, OAD, OAI y \textbf{perfil} (para detectar fugas anteriores o posteriores). Chasis 35$\times$43 apaisado dividido en dos. El órgano reconstruido suele verse dilatado o tortuoso: hallazgo esperable.

\textbf{Protocolo en disfagia progresiva grave:}
\begin{enumerate}[noitemsep]
	\item Iniciar con buche de yodado hidrosoluble bajo radioscopia.
	\item Si no hay fuga ni paso a la vía aérea, proceder con bario.
	\item En estenosis muy severa, \textbf{no administrar sales efervescentes} (riesgo de rotura esofágica por distensión brusca).
\end{enumerate}

% ---
\subsection*{Tránsito del Intestino Delgado}

Consiste en administrar bario oral y documentar su avance por yeyuno e íleon mediante \textbf{radiografías seriadas cada 30--60 minutos}, hasta opacificar el \textbf{ciego} (válvula ileocecal). Es de bajo rendimiento diagnóstico actual; ha sido ampliamente desplazado por la \textbf{enterotomografía} y la \textbf{enterorresonancia}. Aún se utiliza en postoperatorios para evaluar reactivación del tránsito (\textbf{íleo paralítico}).

% ---
\subsection*{Colon por Enema}

Alternativa a la videocolonoscopía (VCC) cuando esta es incompleta por estenosis o colon flexuoso, o en jóvenes con estreñimiento crónico (sospecha de dolicocolon / megacolon). También se usa en el prequirúrgico de reconstrucción de colostomías (ambos cabos con yodado).

\textbf{Preparación:} Dieta + laxantes desde el día anterior (colon libre de materia fecal).

\textbf{Preparación del contraste:} Bario en polvo disuelto en 800 cc -- 1,5 litros de agua tibia; se bate, se presuriza con perita y se clampea la bolsa.

\textbf{Secuencia del estudio:}
\begin{enumerate}[noitemsep]
	\item Paciente en \textbf{decúbito prono}; se introduce cánula por el ano.
	\item \textbf{Fase de llenado:} la bolsa colgada permite entrada de bario por gravedad bajo control radioscopio; cuando llega al colon transverso medio, se baja al piso (maniobra de retorno) para que el exceso regrese a la bolsa dejando la mucosa ``pintada'' hasta el colon derecho.
	\item \textbf{Fase de doble contraste:} se inyecta aire pisando la bolsa manualmente; el aire distiende el colon y el bario adherido revela el relieve mucoso.
	\item \textbf{Documentación:} rotación artesanal del paciente bajo radioscopia para desplegar los \textbf{ángulos hepático y esplénico}. Proyecciones obligatorias: decúbito supino frente, ambas oblicuas, \textbf{perfil de recto} (ampolla rectal y espacio presacro) y \textbf{bipedestación} (niveles hidroaéreos y movilidad).
\end{enumerate}

\textbf{Chasis:} Aproximadamente 5 chasis 35$\times$43 (la mayoría divididos en dos exposiciones).

\textbf{Rol del Licenciado:} Si el médico está a pie de tubo rotando al paciente, el licenciado dispara, centra y cambia chasis rápidamente. Es un estudio artesanal y médico-dependiente.\\

\begin{cajakeypoints}
	\textbf{Principio ``menos es más'':} Usar volumen moderado de bario. Exceso de bario en el colon tapa la mucosa y dificulta ver pólipos o lesiones (imágenes muy blancas con superposición). El Licenciado debe estar con \textbf{chaleco plomado} durante la fase de aire porque opera la bolsa a pie de tubo.
\end{cajakeypoints}

% ============================================================
%  SECCIÓN 4 --- ESTUDIOS DE GLÁNDULAS ANEXAS
% ============================================================
\section*{Estudios de Glándulas Anexas}

\subsection*{Colangiografía Postoperatoria (Tubo en T / Tubo de Kehr)}

Se realiza en el postoperatorio de \textbf{colecistectomía} para evaluar litiasis residual, estenosis postquirúrgica y permeabilidad del contraste al duodeno. Aunque la \textbf{colangiorresonancia} es la técnica de elección actual, este estudio sigue vigente por accesibilidad.

\textbf{Material:} Tubo de Kehr (tubo en T dejado por el cirujano en el colédoco o conducto cístico). Contraste \textbf{yodado hidrosoluble diluido al 50\%} con solución fisiológica (para no tapar pequeños cálculos con una densidad excesiva).

\textbf{Técnica:} Inyección muy lenta y constante con jeringa de 10 cc (evitar espasmos y burbujas de aire, que se confunden con cálculos). Se busca opacificar toda la vía biliar intra y extrahepática.

\textbf{Proyecciones:} Frente en decúbito dorsal + oblicua anterior izquierda (OAI) para desplegar los conductos hepáticos de la columna. Chasis 35$\times$43 dividido en dos.

\textbf{Criterios de evaluación:}
\begin{itemize}[noitemsep]
	\item Ausencia de defectos de relleno (cálculos residuales).
	\item Llenado completo de los conductos hepáticos derecho e izquierdo.
	\item Paso libre del contraste al \textbf{duodeno} (esfínter de Oddi permeable).
\end{itemize}

\subsection*{Sialografía}

Estudio de las glándulas salivales mayores (parótida y submaxilar) para evaluar anatomía de conductos excretores y detectar \textbf{sialolitiasis}. Ha perdido terreno frente a la resonancia, pero persiste para detección de cálculos.

\textbf{Técnica artesanal:}
\begin{enumerate}[noitemsep]
	\item Estimular la secreción con \textbf{jugo de limón} para identificar el orificio ductal.
	\item Cateterizar con \textbf{aguja roma} (sin punta).
	\item Inyectar contraste \textbf{yodado hidrosoluble}.
\end{enumerate}

\textbf{Conductos y localización:}
\begin{itemize}[noitemsep]
	\item \textbf{Conducto de Stenon (parótida):} orificio a nivel del 2\textordfeminine/3\textsuperscript{er} molar superior.
	\item \textbf{Conducto de Wharton (submaxilar):} orificio en el piso de la boca, detrás de la arcada dentaria inferior (carúncula sublingual).
\end{itemize}

\textbf{Proyecciones:} AP y lateral sin contraste (base), luego fase de llenado y fase de vaciado (perfil y oblicuas de mandíbula ``desenfilado'').

% ============================================================
%  SECCIÓN 5 --- HALLAZGOS PATOLÓGICOS RELEVANTES
% ============================================================
\section*{Hallazgos Patológicos Relevantes}

\subsection*{Signos radiológicos: Esófago}

{\renewcommand{\arraystretch}{1.3}
	\begin{center}
		\begin{tabular}{>{\bfseries}p{3.5cm} p{5cm} p{5cm}}
			\hline
			\textbf{Patología} & \textbf{Signos radiológicos} & \textbf{Clínica / Nota} \\
			\hline
			Divertículo de Zenker & Imagen de adición redondeada u ovalada en pared posterior cervical (triángulo de Killian); se rellena de bario antes que el esófago distal. Confirmar en \textbf{perfil estricto}. & Disfagia alta, regurgitación, halitosis. Riesgo: perforación endoscópica si no se detecta antes. \\[0.5em]
			\hline
			Estenosis neoplásica & Defecto de llenado, bordes ``mordidos'' o ``en corazón de manzana'', rigidez segmentaria (sin ondas peristálticas), dilatación suprastenótica + nivel hidroaéreo. & Disfagia progresiva: sólidos $\to$ líquidos $\to$ afagia. \\[0.5em]
			\hline
			Acalasia & ``Punta de lápiz'' / ``cola de ratón'' (signo patognomónico): estrechamiento suave y simétrico en la unión esofagogástrica. Megaesófago / esófago sigmoideo. Ausencia de burbuja gástrica. & Trastorno funcional; esfínter esofágico inferior no se relaja. DD con cáncer: bordes lisos (acalasia) vs. bordes irregulares (neoplasia). \\
			\hline
		\end{tabular}
	\end{center}
}

\subsection*{Signos radiológicos: Estómago}

\textbf{Úlcera gástrica:} El bario rellena el cráter generando una \textbf{imagen de adición} que sobresale del contorno gástrico. Sitio más frecuente: \textbf{curvatura menor}. Benignidad: nicho liso y redondeado con pliegues que convergen simétricamente. Malignidad: nicho dentro de una masa, bordes irregulares, no sobrepasa la línea de la pared.

\textbf{Neoplasia gástrica:} Interfase abrupta e irregular, rigidez segmentaria (zona fija, ``congelada'', sin peristaltismo), defecto de llenado asimétrico.

% ============================================================
%  SECCIÓN 6 --- RADIOPROTECCIÓN
% ============================================================
\section*{Radioprotección en Estudios Contrastados}

Estos estudios implican uso de radiaciones ionizantes. Los efectos biológicos son:
\begin{itemize}[noitemsep]
	\item \textbf{Estocásticos:} aleatorios, proporcionales a la dosis acumulada (ej. riesgo de cáncer).
	\item \textbf{Determinísticos:} aparecen tras superar un umbral de dosis (ej. radiodermitis).
\end{itemize}

\textbf{Principio ALARA:} obtener la mayor calidad de imagen con la menor exposición posible.

\textbf{EPP obligatorio} (el personal suele estar a pie de tubo durante el estudio):

{\renewcommand{\arraystretch}{1.25}
	\begin{center}
		\begin{tabular}{>{\bfseries}p{4cm} p{9cm}}
			\hline
			Elemento & Protege \\
			\hline
			Chaleco plomado & Torso (órganos radiosensibles: tiroides, gónadas, médula ósea). \\
			\hline
			Protector de tiroides & Glándula tiroidea (muy radiosensible). \\
			\hline
			Lentes plomados & Cristalino (riesgo de catarata por dosis acumulada). \\
			\hline
		\end{tabular}
	\end{center}
}
.\\

\begin{cajariesgo}
	\textbf{Importante:} El uso de \textbf{radioscopia} debe minimizarse al máximo (solo para centrar y controlar el paso del contraste). La documentación diagnóstica se realiza con la \textbf{exposición radiográfica en el momento justo}. En el colon por enema, el Licenciado opera la bolsa de aire \textbf{a pie de tubo}: el chaleco es indispensable.
\end{cajariesgo}

% ============================================================
%  ESQUEMA MENTAL --- REPASO FINAL
% ============================================================
\section*{Esquema Mental --- Repaso Final}

\begin{cajakeypoints}
	\begin{center}
		\textbf{\color{azulprincipal}
			Alto Z $\rightarrow$ Efecto fotoeléctrico $\rightarrow$ Mayor densidad en imagen $\rightarrow$ Diferenciación del órgano}\\[0.4em]
		\textbf{\color{naranjadestacado}
			Si hay sospecha de perforación, fuga o fístula: SIEMPRE yodado hidrosoluble, NUNCA bario.}\\[0.3em]
		Bario Z=56 \quad Yodado Z=53 \quad Doble contraste = positivo + negativo\\[0.3em]
		Ayuno EGD: 6 h \quad Orden deglución: semisólido $\to$ líquido $\to$ sólido \quad Colon: laxante día previo\\[0.3em]
		\textbf{Acalasia:} punta de lápiz + megaesófago + sin burbuja gástrica\\[0.3em]
		\textbf{Neoplasia esofágica:} bordes ``corazón de manzana'' + rigidez + dilatación suprastenótica\\[0.3em]
		\textbf{Protocolo EGD:} bipedestación perfil $\to$ OAD bipe $\to$ frente bipe $\to$ decúbito supino $\to$ OAD/OAI $\to$ prono $\to$ OAD duodeno\\[0.3em]
		\textbf{Colangiografía:} Kehr $\to$ yodado 50\% $\to$ inyección lenta $\to$ frente + OAI $\to$ paso libre al duodeno
	\end{center}
\end{cajakeypoints}

\vspace{1em}
\hrule
\vspace{0.5em}
{\footnotesize \textit{Fuente: material de clase --- Dr. N. Asteggiante / Asis. Lic. G. Ferreira,
		Imagenología Especializada I, UdelaR 2026. Bibliografía: Merril Cap. 17; Bontrager Caps. 14--15.}}